% Generated draft: requires mathematical and source review.
\documentclass[11pt]{article}
\usepackage[margin=25mm]{geometry}
\usepackage{fontspec}
\setmainfont{Latin Modern Roman}
\usepackage{amsmath,amssymb,mathtools,unicode-math}
\setmathfont{Latin Modern Math}
\usepackage{enumitem,needspace,xurl,hyperref}
\setlength{\parindent}{0pt}
\setlength{\parskip}{5pt}
\setlength{\emergencystretch}{3em}
\newcommand{\N}{\mathbb{N}}
\newcommand{\Z}{\mathbb{Z}}
\newcommand{\Q}{\mathbb{Q}}
\newcommand{\R}{\mathbb{R}}
\newcommand{\C}{\mathbb{C}}
\newcommand{\F}{\mathbb{F}}
\newcommand{\A}{\mathbb{A}}
\newcommand{\PP}{\mathbb{P}}
\newcommand{\E}{\mathbb{E}}
\newcommand{\eps}{\varepsilon}
\newcommand{\dd}{\,\mathrm{d}}
\begin{document}
\section*{UnsolvedMath ID: 30005850}
\textbf{Source problem code:} catalogue entry
\subsection{Definitions and mathematical statement}
The supplied record identifies the setting as the Dong--Song positive-mass stability theorem, referred to as Theorem 4 in the source discussion. Let \(Z_i\) denote the sets called the \emph{bad sets} in that theorem, indexed by the members of the sequence under consideration. The theorem itself, its hypotheses, the ambient spaces containing \(Z_i\), and the precise definition of these sets are not supplied. The terms \emph{splines}, \emph{bubbles}, and \emph{drawstrings} name geometric phenomena in the source discussion, but their mathematical definitions are also absent.

The source question is therefore retained in the following conditional form: under the full hypotheses and with the bad sets \(Z_i\) of the referenced Dong--Song theorem, do splines, bubbles, and drawstrings exhaust the phenomena that can occur on \(Z_i\)? The expression ``exhaust the phenomena'' requires a precise classification criterion; the supplied record does not specify whether it means a decomposition of each \(Z_i\), a description of limiting geometry, or another notion. In particular, no disjoint-union assertion or exhaustive quantified classification is justified by the supplied text.

The source proposes two complementary directions for obtaining evidence:
\begin{enumerate}
\item Formulate and prove a stability conjecture with an additional condition ruling out splines and bubbles, and with a conclusion of \(d_p\) stability.
\item Formulate and prove another stability conjecture with an additional condition ruling out bubbles and drawstrings, and with a conclusion of volume-preserving Sormani--Wenger intrinsic flat convergence.
\end{enumerate}
The additional conditions, the definition and parameter range of \(d_p\), the convergence targets, and the remaining hypotheses are not specified. These are proposed formulations to be developed, not stated implications with established hypotheses. The source describes these two proposed results, together with the referenced stability result, as potential strong evidence for the classification, rather than as a proof that it is exhaustive.

The broader motivation concerns sequences with scalar-curvature lower bounds. The source explicitly qualifies its proposed list by observing that \emph{sewing} can also occur if the topology of the sequence is allowed to vary. No fixed-topology hypothesis or precise definition of sewing is supplied, so the three named phenomena cannot be asserted to give an exhaustive classification for arbitrary topology-varying sequences.

The excerpt also includes a separately numbered Question 2 asking whether one can formulate and prove a version of Geroch stability in the style of Theorem 4. Its inclusion in the present record is unexplained; it is retained as an adjacent question, without identifying it with the classification question. A fully precise mathematical formulation requires the missing theorem, definitions, hypotheses, and intended classification criterion.
\subsection{Short English statement}
The question asks whether the bad sets in the referenced Dong--Song positive-mass stability theorem exhibit only splines, bubbles, and drawstrings. Two further stability conjectures are suggested as evidence: excluding splines and bubbles should lead to \(d_p\) stability, while excluding bubbles and drawstrings should lead to volume-preserving Sormani--Wenger intrinsic flat convergence. The source warns that sewing is also possible when topology varies. Missing definitions and theorem hypotheses prevent a fully quantified statement.
\subsection{Sources}
\begin{itemize}
\item Local UnsolvedMath registry, permanent ID \(30005850\), title \emph{Classification of Bad Sets in Scalar-Curvature Limits}, exact source code \texttt{catalogue entry}. The record attributes the problem to a 2024 Oberwolfach workshop discussion; no individual proposer is supplied.
\item \emph{Analysis, Geometry and Topology of Positive Scalar Curvature Metrics} (2024), Oberwolfach Reports, pp.\ 483--566, DOI \texttt{10.4171/owr/2024/9}: \url{https://doi.org/10.4171/owr/2024/9}. The supplied excerpt contains Question 3 and an appended Question 2, and refers to Theorem 4 without reproducing it.
\item Supplied literature source: \emph{Stability of Llarull's theorem in all dimensions}, \url{https://www.sciencedirect.com/science/article/abs/pii/S0001870824004961}.
\item Supplied literature source: \emph{Conjectures on convergence and scalar curvature}, \url{https://arxiv.org/abs/2103.10093}.
\item Literature assessment dated exactly \texttt{2026-08-22}: the supplied curation labels the problem \emph{partially solved} and reports that recent Llarull and scalar-curvature stability theorems control several bad-set phenomena, including splines in important settings, while no theorem proves that all bad sets consist only of splines, bubbles, and drawstrings. This is supplied dated curation, not an independently verified literature assessment.
\item The registry labels the statement \texttt{corrected\_verified} and reports reconstruction and checking against the official Oberwolfach report. These are supplied source-verification claims; no independent source visit or verification is asserted here.
\end{itemize}
\end{document}
